%&pdftex
%%%%%%%%%%%%%%%%%%%%%%%%%%%%%%%%%%%%%%%%%%%%%%%%%%%%%%%%%%%%%%%%%%%%%%%%%%%%%%%
%                                                                             %
%           ON COMPLETENESS OF ORTHOGONAL SYSTEMS AND DIRAC DELTAS,           %
%                                                                             %
%                     POR MARIO PEREZ Y JUAN LUIS VARONA.                     %
%                                                                             %
%     PENULTIMA MODIFICACION: 2 DE JUNIO DE 1993, POR MARIO PEREZ RIERA.      %
%       ULTIMA MODIFICACION: 22 DE OCTUBRE DE 1993, POR JUAN LUIS VARONA.     %
%                                                                             %
%                      ENVIADO PARA SU PUBLICACION EN                         %
%              Journal of Computational and Applied Mathematics               %
%                                                                             %
%%%%%%%%%%%%%%%%%%%%%%%%%%%%%%%%%%%%%%%%%%%%%%%%%%%%%%%%%%%%%%%%%%%%%%%%%%%%%%%

\magnification = 1200

%\baselineskip=24pt  % POR NORMAS DE LA REVISTA UNICAMENTE.

\font \peq = cmr8
\font \may = cmcsc10

\font\msbmnormal=msbm10
\font\msbmpeq=msbm7
\font\msbmmuypeq=msbm5
\newfam\numeros
\textfont\numeros=\msbmnormal
\scriptfont\numeros=\msbmpeq
\scriptscriptfont\numeros=\msbmmuypeq
\def\num{\fam\numeros\msbmnormal}
\def\R{{\num R}}
\def\C{{\num C}}
\def\T{{\num T}}
\def\Z{{\num Z}}

\def \proof {{\it Proof.} }
\def \qed {\vrule height 5 pt width 5 pt depth 0 pt}

\centerline{\bf ON COMPLETENESS OF ORTHOGONAL SYSTEMS}
\centerline{\bf AND DIRAC DELTAS}

\bigskip

\centerline{\may Mario P\'erez Riera and Juan L. Varona Malumbres}

\footnote{}{1991 Mathematics Subject Classification: 42C05, 42C30.}
\footnote{}{Keywords: orthogonal systems, Dirac deltas, moment problem.}
\footnote{}{Research supported by DGICYT PB89-0181-C02-02 and by IER.}
\footnote{}{PAPER PUBLISHED IN:
{\it J. Comput.\ Appl.\ Math.}\ {\bf 58} (1995), 225--231.}
\bigskip

{\narrower\narrower\noindent{\sl Given a positive measure $\mu$ supported on a
set $\Omega \subseteq \C$, an  orthonormal system $\{\varphi_n\}_{n\geq 0}$ and
a point $a \in \Omega$, we study the relationship among $\mu(\{a\})$, the
kernels $K_n(a,a) = \sum_{k=0}^n \varphi_k(a) \overline{\varphi_k(a)}$ and the
denseness of ${\rm span}\{\varphi_n\}_{n\geq 0}$ in $L^2(\mu)$ and in
$L^2(\nu)$, where $\nu = \mu + M \delta_a$.}\smallskip}

\bigskip
\centerline {\may 0. Introduction}
\bigskip

Let $\mu$ be a positive measure supported on a subset $\Omega \subseteq \C$ and
$\{\varphi_n \colon \Omega \to \C\}_{n\geq 0}$ an orthonormal system in
$L^2(\mu)=L^2(\Omega, \mu)$. Then
$$
   \int_{\Omega}\varphi_n \overline \varphi_m\,d\mu =
   \cases{0,&if $n \not= m$; \cr 1, &if $n=m$.}
$$
The system $\{\varphi_n\}_{n\geq 0}$ is said to be complete in $L^2(\mu)$ if
the set ${\rm span}\{\varphi_n\}_{n\geq 0}$ of finite linear combinations
is dense in $L^2(\mu)$ or, in other terms, if for each $\Phi \in L^2(\mu)$
$$
   \int_{\Omega}\Phi \overline\varphi_n \,d\mu = 0 \quad \forall n \geq 0
   \iff \Phi = 0 \quad \mu\hbox{-a.e.}
$$
(the orthogonality is not required here).

For each $n \geq 0$, set $\Pi_n = \Big\{\displaystyle \sum_{k=0}^n \lambda_k
\varphi_k; \lambda_0, \dots, \lambda_n \in \C\Big\}$. The best $L^2(\mu)$
approximant in $\Pi_n$ of any $f \in L^2(\mu)$ is given by the $n$-th partial
sum of its Fourier series with respect to the set $\{\varphi_n\}_{n \geq 0}$.
Thus the best approximant is
$$
   S_n(f,x) = \sum_{k=0}^n c_k(f) \varphi_k(x) =
   \int_{\Omega} f(y) K_n(x,y)\,d\mu(y),
$$
where
$$
   c_k(f) = \int_{\Omega} f \overline\varphi_k \,d\mu, \quad
   K_n(x,y) = \sum_{k=0}^n \varphi_k(x) \overline {\varphi_k(y)}.
$$
Furthermore,
$$
   \sum_{k=0}^{\infty} |c_k(f)|^2 \leq
   \int_{\Omega} |f|^2 \,d\mu \quad \forall f \in L^2(\mu)
   \eqno (1)
$$
(Bessel's inequality). If the system $\{\varphi_n\}_{n \geq 0}$ is complete in
$L^2(\mu)$, then (1) becomes an equality (Parseval's equality) and the
approximants $\{S_n(f)\}_{n \geq 0}$ converge to $f$ in $L^2(\mu)$. This
leads to an elementary proof of the following result (see [1, p.~63, 114], [6,
p.~45] for more elaborate proofs and only in the case of systems of
polynomials):

\proclaim{\may Proposition 1}. Let $\{\varphi_n\}_{n \geq 0}$ be
an orthonormal system in $L^2(\mu)$ and let $a \in \Omega$. Then
$$
   {1 \over K_n(a,a)} \geq \mu(\{a\}) \quad \forall n \geq 0.
   \eqno (2)
$$
If $\mu(\{a\}) > 0$ and $\{\varphi_n\}_{n \geq 0}$ is complete, then
$$
   \lim_n {1 \over K_n(a,a)} = \mu(\{a\}).
   \eqno (3)
$$

\proof We can assume $0 < \mu(\{a\}) < \infty$, otherwise the statement holds
trivially. Let $f$ be the characteristic function at the point $a$; then
$c_k(f) = \overline {\varphi_k(a)} \mu(\{a\})$, so that
$$
   \sum_{k=0}^{\infty} |c_k(f)|^2 =
   \mu(\{a\})^2 \sum_{k=0}^{\infty} |\varphi_k(a)|^2
$$
and
$$
   \int_{\Omega} |f|^2 \,d\mu = \mu(\{a\}).
$$
Now, (2) and (3) follow from Bessel's and Parseval's formula, respectively.
\qed

\medskip

Concerning the measure $\mu$ and the kernels $K_n(a,a)$, we have the  
following well-known, elementary result (see [3, p.~38, theorem 7.3] or [4,
p.~4], for example). In fact, inequality (2) can also be obtained as a
corollary of lemma 1.

\proclaim {\may Lemma 1}. Let $\{\varphi_n\}_{n \geq 0}$ be an orthonormal
system in $L^2(\mu)$ and let $a \in \Omega$. Then
$$
   {1 \over K_n(a,a)} = \min \int_{\Omega} |R_n|^2 \,d\mu,
$$
where the minimum is taken over all $R_n \in \Pi_n$ such that $R_n(a) = 1$.
Furthermore, this minimum is attained for $R_n(x) = K_n(x,a)/K_n(a,a)$.

Our aim is to use proposition 1 and lemma 1 to obtain, using elementary
techniques, some relations between $\displaystyle \lim_n K_n(a,a)^{-1}$ and
certain properties of completeness of the system $\{\varphi_n\}_{n\geq0}$.

\bigskip
\centerline {\may 1. Addition of a mass point}
\bigskip

Obviously, one can not expect (3) to hold if the system $\{\varphi_n\}_{n\geq
0}$ is not complete. If $\{\varphi_n\}_{n\geq 0}$ is complete but $\mu(\{a\}) =
0$, it can also fail, for the values $\varphi_n(a)$ are $\mu$-meaningless; in
this case, we will prove that (3) holds if and only if the system
$\{\varphi_n\}_{n\geq 0}$ is complete in $L^2(\nu)$. Here $\nu = \mu +
M\delta_a$, where $\delta_a$ is a Dirac delta on $a$, $M>0$ and as a consequence
$$
   \int_{\Omega} f \,d\nu = \int_{\Omega} f \,d\mu + M f(a).
$$

The system $\{\varphi_n\}_{n\geq 0}$ may not be orthogonal in $L^2(\nu)$, but
it can be orthonormalized so as to get an orthonormal system $\{\psi_n\}_{n\geq
0}$ in $L^2(\nu)$, such that $\psi_n = \displaystyle \sum_{k=0}^n
\lambda_{n,k} \varphi_k$, with $\lambda_{n,n} \not= 0$. Clearly, if
$\{\psi_n\}_{n\geq 0}$ (or, equivalently, $\{\varphi_n\}_{n\geq 0}$) is
complete in $L^2(\nu)$, then $\{\varphi_n\}_{n\geq 0}$ is also complete in
$L^2(\mu)$, but the converse is not true, in general.

In view of (3), we will mainly deal with the case $\mu(\{a\}) = 0$. However,
note that if $\mu(\{a\}) > 0$ then the measures $\nu$ and $\mu$ are equivalent
and so $\{\varphi_n\}_{n\geq 0}$ is complete in $L^2(\mu)$ if and only if
$\{\psi_n\}_{n\geq 0}$ is complete in $L^2(\nu)$ (for the same reason, $M$
could be taken equal to 1).

Let us state our first result (another proof of part ``$a) \Longrightarrow b)$''
can be found in [2, lemma 2]):

\proclaim {\may Theorem 1}. If $\{\varphi_n\}_{n\geq 0}$ is a complete
orthonormal system in $L^2(\mu)$, $\mu(\{a\}) = 0$ and $\nu = \mu + M \delta_a$,
then the following properties are equivalent:
\item{a)} $\displaystyle \lim_n {1 \over K_n(a,a)} = 0$;
\item{b)} $\{\psi_n\}_{n\geq 0}$ is a complete orthonormal system in $L^2(\nu)$.

\proof $a) \Longrightarrow b)$: Suppose $\{\psi_n\}_{n\geq 0}$ is not complete
in $L^2(\nu)$; then, there exists $\Phi \in L^2(\nu)$, $\Phi \not= 0$, such
that
$$
   \int_{\Omega} \Phi \overline \psi_n \,d\nu = 0 \quad \forall n \geq 0.
$$
We can also assume
$$
   \int_{\Omega} |\Phi|^2 \,d\nu = 1,
$$
so that $\{\Phi\}\cup\{\psi_n\}_{n\geq 0}$ is an orthonormal system in
$L^2(\nu)$. Furthermore, $\Phi(a) \not= 0$, otherwise it would be orthogonal to
$\{\varphi_n\}_{n\geq 0}$ in $L^2(\mu)$ and therefore $\Phi = 0$ $\mu$ and
$\nu$ a.e.

Put $D_n(a,a) = \displaystyle \sum_{k=0}^n |\psi_k(a)|^2$. Then, by (2)
applied to $\{\Phi\}\cup\{\psi_n\}_{n\geq0}$, the fact $\Phi(a)\not=0$, and
lemma 1 respectively, we have the chain of inequalities
$$
   M = \nu(\{a\}) \leq \lim_n {1 \over |\Phi(a)|^2 + D_n(a,a)} <
   \lim_n {1 \over D_n(a,a)} = \lim_n {1 \over K_n(a,a)} + M
$$
and so
$$
   \lim_n {1 \over K_n(a,a)} > 0,
$$
which is a contradiction.

$b) \Longrightarrow a)$: By (3) applied to $\{\Phi\}\cup\{\psi_n\}_{n\geq0}$
and lemma 1,
$$
   M = \lim_n {1 \over D_n(a,a)} = \lim_n {1 \over K_n(a,a)} + M
$$
which gives a). \qed

\medskip

Under the conditions of theorem 1, the orthonormal system $\{\psi_n\}_{n\geq0}$
may not be complete in $L^2(\nu)$, but in this case it becomes complete by
adding just one new function:

\proclaim {\may Proposition 2}. Let $\{\varphi_n\}_{n\geq 0}$ be a complete
system in $L^2(\mu)$, $\mu(\{a\}) = 0$ and $\nu = \mu + M \delta_a$ and suppose
$\{\psi_n\}_{n\geq 0}$ is not complete in $L^2(\nu)$. Then, the system
$\{\Phi\}\cup\{\psi_n\}_{n\geq0}$ is orthogonal ($\Phi$ is not normalized) and
complete in $L^2(\nu)$, where
$$
   \Phi(x) = \cases{
   \displaystyle \sum_{k=0}^{\infty} \varphi_k(x) \overline {\varphi_k(a)},
   &if $x \not= a$;
   \cr
   \cr -{1 \over M}, &if $x = a$.}
$$

\proof By theorem 1,
$$
   \sum_{k=0}^{\infty} |\varphi_k(a)|^2 < \infty.
$$
Then, as
$$
    \int_{\Omega} \Bigl|\sum_{k=n}^m
        \varphi_k(x) \overline{\varphi_k(a)}\Bigr|^2 \,d\nu(x)
$$
$$
    = M \sum_{k=n}^m |\varphi_k(a)|^2
    + \int_{\Omega} \Bigl|\sum_{k=n}^m
        \varphi_k(x) \overline{\varphi_k(a)}\Bigr|^2 \,d\mu(x)
    = (M+1) \sum_{k=n}^m |\varphi_k(a)|^2,
$$
the series $\displaystyle \sum_{k=0}^{\infty} \varphi_k(x)
\overline{\varphi_k(a)}$ converges in $L^2(\nu)$ because its partial sums
constitute a Cauchy sequence in $L^2(\nu)$. So $\Phi$ is a well defined
function in $L^2(\nu)$. Now,
$$
   \int_{\Omega} \varphi_n \overline \Phi \,d\nu =
   M \varphi_n(a) \overline {\Phi(a)} +
   \sum_{k=0}^{\infty} \varphi_k(a)
   \int_{\Omega} \varphi_n \overline \varphi_k \,d\mu = 0
$$
for every $n \geq 0$; therefore, we also have
$$
   \int_{\Omega} \psi_n \overline \Phi \,d\nu = 0 \quad \forall n \geq 0
$$
and $\{\Phi\} \cup \{\psi_n\}_{n\geq 0}$ is an orthogonal system in $L^2(\nu)$.
In order to prove that it is complete in $L^2(\nu)$, it is enough to check that
$$
   f \in L^2(\nu), \quad \int_{\Omega} f \overline\psi_n \,d\nu = 0
   \quad \forall n \geq 0 \quad \Longrightarrow\quad f = C \Phi \quad
   \nu\hbox{-a.e.}
$$
If
$$
   \int_{\Omega} f \overline\psi_n \,d\nu = 0 \quad \forall n \geq 0
$$
then
$$ 
   \int_{\Omega} f \overline\varphi_n \,d\nu = 0 \quad \forall n \geq 0.
$$
Thus,
$$
   \int_{\Omega} (f+Mf(a)\Phi) \overline\varphi_n \,d\mu =
   \int_{\Omega} (f+Mf(a)\Phi) \overline\varphi_n \,d\nu = 0
   \quad \forall n \geq 0,
$$
i.e.,
$$
   f+Mf(a)\Phi = 0 \quad \mu\hbox{-a.e.}
$$
and so
$$
   f+Mf(a)\Phi = 0 \quad \nu\hbox{-a.e.}\ \qed
$$

\medskip

\noindent{\may Example.} Let $\T$ be the unit circle, $\Omega = \T \cup \{0\}$
and $\mu$ the measure
$$
   \int_{\Omega} f \,d\mu
   = {1 \over 2\pi}\int_{-\pi}^{\pi} f(e^{i\theta}) \,d\theta.
$$
Let $\{\varphi_n\}_{n\in\Z}$ be given by $\varphi_n(z) = z^n$ for $n \geq 0$
and $\varphi_n(z) = ({\overline z})^{-n}$ for $n < 0$ (the system is indexed in
$\Z$, but this makes no difference). This system is orthonormal and complete in
$L^2(\mu)$. Take $\nu = \mu + \delta_0$. Then, the system
$\{\varphi_n\}_{n\in\Z}$ is not complete in $L^2(\nu)$ and $\Phi(0) = -1$,
$\Phi(z) = 1$ for $z \in \T$.

\medskip

Given any finite positive measure $\nu$ supported on the unit circle, the
orthonormal system obtained from $\{z^n\}_{n\in\Z}$ is complete in $L^2(\nu)$
(see, e.g. [1, p.~180, theorem 5.1.2]). Then, if we consider a finite positive
measure $\mu$ supported on the unit circle and the orthonormal system
$\{\varphi_n\}_{n\in\Z}$ obtained from $\{z^n\}_{n\in\Z}$, part b) in theorem
1 holds, so that part a) also holds.

In contrast, the real case is more interesting, as we see in the next section.

\bigskip
\centerline {\may 2. Orthogonal polynomials on the real line}
\bigskip

In the following we will consider a system $\{p_n\}_{n\geq 0}$ of polynomials
($p_n$ of degree $n$) orthonormal with respect to some positive measure $\mu$
on $\R$.

The measure $\mu$ is said to be determinate if there does not exist any other
positive measure $\eta$ on $\R$ such that
$$
   \int_{\R} x^n \,d\mu(x) = \int_{\R} x^n \,d\eta(x) \qquad \forall n \geq 0;
$$
otherwise, $\mu$ is said to be indeterminate.

The system $\{p_n\}_{n\geq 0}$ is complete in $L^2(d\mu)$ if and only if $\mu$
is {\it $N$-extremal} (see [5]). Every determinate measure is $N$-extremal, and
every indeterminate $N$-extremal measure is a countable sum of Dirac deltas
(see [2]).

If $\mu$ is determinate, then
$$
   \lim_n {1 \over K_n(a,a)} = \mu(\{a\}) \qquad \forall a \in \R \eqno (4)
$$
(see, e.g., [6, p.~45, corollary 2.6]), while if $\mu$ is
indeterminate then
$$
   \lim_n {1 \over K_n(a,a)} > 0 \qquad \forall a \in \C \eqno (5)
$$
(see, e.g., [1, p.~50] and [6, p.~50, corollary 2.7]).

Now, the previous results provide a simple proof of the following:

\proclaim{\may Theorem 2}. Let $\{p_n\}_{n\geq 0}$ be a system of polynomials
orthonormal with respect to a positive measure $\mu$ on $\R$. Let $a \in \R$,
$M > 0$, $\nu = \mu + M\delta_a$. Then:
\item{a)} $\mu$ indeterminate $N$-extremal, $\mu(\{a\}) = 0$ $\Longrightarrow$
$\nu$ indeterminate not $N$-extremal.
\item{b)} $\mu$ indeterminate $N$-extremal, $\mu(\{a\}) > 0$ $\Longrightarrow$
$\nu$ indeterminate $N$-extremal.
\item{c)} $\mu$ indeterminate not $N$-extremal $\Longrightarrow$ $\nu$
indeterminate not $N$-extremal.
\item{d)} $\mu$ determinate, $\mu(\{a\}) = 0$
$\Longrightarrow$ $\nu$ determinate or indeterminate $N$-extremal.
\item{e)} $\mu$ determinate, $\mu(\{a\}) > 0$ $\Longrightarrow$ $\nu$
determinate.

\proof $a)$, $b)$ and $c)$: since $\mu$ is indeterminate, $\nu = \mu +
M \delta_a$ is also indeterminate. Now, if $\mu$ is not $N$-extremal (i.e. the
polynomials are not dense in $L^2(\mu)$) then clearly $\nu$ is not
$N$-extremal. If $\mu$ is $N$-extremal and $\mu(\{a\}) > 0$, then $\nu$ is also
$N$-extremal, for both measures are equivalent. Finally, if $\mu(\{a\}) = 0$,
from (5) and theorem 1 it follows that the polynomials are not dense in
$L^2(\nu)$.

$d)$ and $e)$: if $\mu$ is determinate, from (4) and theorem 1 it follows that
the polynomials are dense in $L^2(\nu)$, so that $\nu$ is either determinate or
indeterminate $N$-extremal. This proves d). Now, assume $\mu(\{a\}) > 0$. Then,
$\mu$ and $\nu$ are equivalent measures. Take $b \in \R$ such that $\mu(\{b\}) =
0$. Applying d), the measure $\mu + \delta_b$ is $N$-extremal. Since $\mu +
\delta_b$ is equivalent to $\nu + \delta_b$, this measure is also $N$-extremal.
By part a), $\nu$ can not be indeterminate $N$-extremal, so that it is
determinate. \qed

\medskip

\noindent{\may Remark.} Both cases in part d) can actually occur. Indeed,
if
$$
   \nu = \sum_{k=0}^{\infty} M_k \delta_{a_k}
$$
is either indeterminate $N$-extremal (every indeterminate $N$-extremal measure
is of this form) or determinate (take, for example, $\{a_k\}_{k\geq 0}$
bounded), it can be shown that the measure
$$
   \mu = \sum_{k=1}^{\infty} M_k \delta_{a_k}
$$
is determinate. A proof can be seen in [2]; it is also a consequence of
inequality (5) and theorem 1.

In this context, let us mention that, in case d), if the measure $\mu$ is not
discrete, then $\nu$ is not discrete; therefore, it must be also determinate.

\bigskip

\centerline {\may References}

\bigskip
\frenchspacing

\item {{\bf [1]}} N. I. Akhiezer,
{\it The Classical Moment Problem\/}
(Oliver and Boyd, Edinburgh, 1965).

\item {{\bf [2]}} C. Berg and J. P. R. Christensen,
Density questions in the classical theory of moments,
{\it Ann. Inst. Fourier} {\bf 31} (1981), 99--114.

\item {{\bf [3]}} T. S. Chihara,
{\it An Introduction to Orthogonal Polynomials\/}
(Gordon and Breach, New York, 1978).

\item {{\bf [4]}} P. Nevai,
Orthogonal Polynomials,
{\it Mem. Amer. Math. Soc.} {\bf 213} (1979).

\item {{\bf [5]}} M. Riesz,
Sur le probl\`eme des moments et le th\'eor\`eme de Parseval correspondant,
{\it  Acta Litt. Ac. Sci. Szeged} {\bf 1} (1923), 209--225.

\item {{\bf [6]}} J. A. Shohat and J. D. Tamarkin,
{\it The Problem of Moments\/}
(Math. Surveys No.~1, Amer. Math. Soc., Providence, RI, 1970).

\bigskip

\halign{# \hskip 3cm\hfill &#\hfill \cr
  Mario P\'erez Riera          &Juan L. Varona Malumbres\cr
  Dpto. de Matem\'aticas       &Dpto. de Matem\'aticas y Computaci\'on\cr
  Universidad de Zaragoza      &Universidad de La Rioja\cr
  50009 Zaragoza               &26004 Logro\~no\cr
  Spain.                       &Spain.\cr}

\bigskip

\centerline{e-mail: jvarona@cc.unizar.es}

\bye